% !TeX root = ../main.tex
\section{Introduction}
The first ICWSM, in 2007, studied weblogs. Blogs appear in 35.2\% of contributions from the conference's first five editions and in 0.9\% of contributions from its most recent five. By 2015, the conference had dropped ``Weblogs'' from its name while keeping its acronym. Twitter then rose to 45.7\% of contributions in 2012--2016 and fell to 24.8\% in 2022--2026. Each shift raises the same question: does research on an earlier platform still describe how people interact online? As AI systems increasingly rank, rewrite, generate, and moderate online interaction \citep{hancock2020,hohenstein2023}, the field faces a sharper version of that question: is social media still the right object of study, or have its central questions moved to AI?

ICWSM's twentieth edition offers a record of what survives these shifts. We separate the subjects the conference studied from the problems researchers asked about them, then trace methodological questions across platform contexts.

Topics alone cannot capture this history. Researchers can study the same subject while pursuing different questions. Community research can examine how participation develops, how members provide support, or how rules are enforced. Methods also shape what can be learned. Studies have evaluated Twitter sampling \citep{p14401}, tested adjustments for textual confounding \citep{p19362}, and used randomized prompts or quasi-random server assignment to study communication \citep{p19308,p31365}. Following subjects, problem framings, and forms of evidence together gives a fuller account of the field.

We build on computational histories of research communities \citep{hall2008,liu2014,wallace2017} and the reflective approach of the FAccT retrospective \citep{facct2022}. Within ICWSM, work on geographical representation shows why diversity of subjects or platforms cannot stand in for diversity of populations studied \citep{weird2024}. We ask:
\begin{itemize}
\item \textbf{RQ1:} How have ICWSM's topics and platform contexts evolved across editions?
\item \textbf{RQ2:} How have problem framings changed across and within research topics?
\item \textbf{RQ3:} How do selected studies illustrate methodological developments and recurring challenges across changing platform contexts?
\end{itemize}

We analyze titles and abstracts from 2,139 indexed contributions across 2007--2026. We distinguish topics, or the subjects and tasks represented in scholarly text, from problem framings, or the questions and concerns expressed within those subjects. Topic modeling describes the former; framing cues approximate the latter. Four fixed five-edition periods provide a common basis for comparison. Selected literature supplies context for interpreting methods without estimating their adoption across the conference.

The results show topical continuity alongside changing problem framing. The community component remains near one tenth of the topic mixture, while governance-related framing becomes much more common. The contrast persists when we remove governance features from the topic representation and restrict abstracts to a fixed length. Harm-related cues also rise under fixed-length checks.

We contribute a documented corpus of twenty ICWSM editions with longitudinal topic, framing, and platform measures. The analysis separates topic change from changes in problem framing, with community research and governance as a detailed case. We also synthesize recurring methodological problems across platform contexts and turn them into a reporting checklist for AI-mediated interaction. The corpus and analysis materials will be released upon publication. Platforms change; the questions recur.
